\documentclass[11pt,a4paper]{article}
\usepackage[utf8]{inputenc}
\usepackage[T1]{fontenc}
\usepackage[portuguese]{babel}
\usepackage{geometry}
\geometry{margin=2.5cm}
\usepackage{listings}
\usepackage{xcolor}
\usepackage{amsmath}
\usepackage{amssymb}
\usepackage{hyperref}
\usepackage{enumitem}
\usepackage{tcolorbox}
\usepackage{tabularx}
\usepackage{booktabs}

\definecolor{codebg}{rgb}{0.95,0.95,0.95}
\definecolor{codegreen}{rgb}{0.1,0.5,0.1}
\definecolor{codepurple}{rgb}{0.5,0.1,0.5}
\definecolor{codegray}{rgb}{0.5,0.5,0.5}

\lstdefinestyle{python}{
    backgroundcolor=\color{codebg},
    commentstyle=\color{codegray}\itshape,
    keywordstyle=\color{codepurple}\bfseries,
    stringstyle=\color{codegreen},
    basicstyle=\ttfamily\small,
    breaklines=true,
    frame=single,
    language=Python,
    showstringspaces=false,
    tabsize=4,
    literate={á}{{\'a}}1 {ã}{{\~a}}1 {é}{{\'e}}1 {í}{{\'i}}1 {ó}{{\'o}}1 {õ}{{\~o}}1 {ú}{{\'u}}1 {ç}{{\c{c}}}1 {Á}{{\'A}}1 {Ã}{{\~A}}1 {É}{{\'E}}1 {Í}{{\'I}}1 {Ó}{{\'O}}1 {Õ}{{\~O}}1 {Ú}{{\'U}}1 {Ç}{{\c{C}}}1 {â}{{\^a}}1 {ê}{{\^e}}1 {ô}{{\^o}}1 {û}{{\^u}}1 {Â}{{\^A}}1 {Ê}{{\^E}}1 {Ô}{{\^O}}1 {Û}{{\^U}}1 {à}{{\`a}}1 {è}{{\`e}}1 {ì}{{\`i}}1 {ò}{{\`o}}1 {ù}{{\`u}}1
}
\lstset{style=python}

\newtcolorbox{objetivos}[1][]{
  colback=green!5!white,
  colframe=green!40!black,
  title=O que você vai aprender nesta página,
  fonttitle=\bfseries,
  #1
}

\title{\textbf{Data Analysis with Python 2024--2025}\\Parte 5: Machine Learning (Regressão Linear)\\\large Conteúdo Teórico}
\author{Tradução e Adaptação: University of Helsinki --- MOOC.fi}
\date{}

\begin{document}
\maketitle

\section*{Nota Importante}
Este material é uma tradução e adaptação baseada no curso \textit{Data Analysis with Python 2024-2025}, oferecido pela \textbf{The University of Helsinki MOOC Center}.

\begin{objetivos}
\begin{itemize}
    \item Compreender o conceito de regressão linear.
\end{itemize}
\end{objetivos}

\tableofcontents
\newpage

\section{Regressão Linear}

A análise de regressão tenta explicar as relações entre variáveis. Uma dessas variáveis, chamada de variável dependente, é o que queremos "explicar" usando uma ou mais variáveis explicativas (ou independentes). Na regressão linear, assumimos que a variável dependente pode ser, aproximadamente, expressa como uma combinação linear das variáveis explicativas. 

Como um exemplo simples, podemos ter a variável dependente "altura" e uma variável explicativa "idade". A idade de uma pessoa pode explicar muito bem a sua altura, e essa relação é aproximadamente linear para crianças (idades entre 1 e 16 anos). Outra forma de pensar sobre a regressão é ajustar uma curva aos pontos de dados observados. Se temos apenas uma variável explicativa, isso é fácil de visualizar, como veremos a seguir.

Podemos aplicar a regressão linear facilmente com o pacote \texttt{scikit-learn}. Vamos analisar alguns exemplos.

Primeiro, fazemos as importações padrão.
\begin{lstlisting}
import numpy as np
import matplotlib.pyplot as plt
import sklearn # Isso importa a biblioteca scikit-learn
\end{lstlisting}

Em seguida, criamos alguns dados com aproximadamente a relação $y = 2x + 1$, com erros normalmente distribuídos.
\begin{lstlisting}
np.random.seed(0)
n = 20 # Numero de pontos de dados
x = np.linspace(0, 10, n)
y = x * 2 + 1 + 1 * np.random.randn(n) # Desvio padrao 1
print(x)
print(y)
\end{lstlisting}

O próximo passo é importar a classe \texttt{LinearRegression}.
\begin{lstlisting}
from sklearn.linear_model import LinearRegression
\end{lstlisting}

Agora podemos ajustar uma linha através dos pontos de dados $(x, y)$:
\begin{lstlisting}
model = LinearRegression(fit_intercept=True)
model.fit(x[:, np.newaxis], y)
xfit = np.linspace(0, 10, 100)
yfit = model.predict(xfit[:, np.newaxis])
plt.plot(xfit, yfit, color="black")
plt.plot(x, y, 'o')
# O codigo a seguir desenhara tantos segmentos de reta quantas forem as colunas nas matrizes x e y
plt.plot(np.vstack((x, x)), np.vstack([y, model.predict(x[:, np.newaxis])]), color="red")
\end{lstlisting}

A regressão linear tenta minimizar a soma dos erros quadráticos:
\[ \sum_{i}(y[i]-\hat{y}[i])^{2} \]
que equivale à soma dos comprimentos ao quadrado dos segmentos de reta (que no exemplo estariam em vermelho) no gráfico traçado. Os valores estimados $\hat{y}[i]$ são denotados por \texttt{yfit[i]} no código acima.

\begin{lstlisting}
print("Parameters:", model.coef_, model.intercept_)
print("Coefficient:", model.coef_[0])
print("Intercept:", model.intercept_)
\end{lstlisting}

Neste caso, o \textit{coefficient} (coeficiente) é a inclinação da linha ajustada, e o \textit{intercept} (intercepto) é o ponto onde a linha ajustada intercepta o eixo y.

Os parâmetros estimados pelo algoritmo de regressão costumam ser bem próximos dos parâmetros que geraram os dados (neste caso, coeficiente 2 e intercepto 1). 

\section{Múltiplas Features (Variáveis Explicativas)}

O exemplo anterior tinha apenas uma variável explicativa. Algumas vezes, isso é chamado de regressão linear simples. O próximo exemplo ilustra uma regressão mais complexa com múltiplas variáveis explicativas.

\begin{lstlisting}
sample1 = np.array([1, 2, 3]) # As tres variaveis explicativas tem valores 1, 2 e 3, respectivamente
sample2 = np.array([4, 5, 6]) # Outro exemplo de valores das variaveis explicativas
sample3 = np.array([7, 8, 10]) # ...

# Para os valores 1, 2 e 3 das variaveis explicativas, o valor y=15 foi observado, e assim por diante.
y = np.array([15, 39, 66]) + np.random.randn(3)
\end{lstlisting}

Vamos tentar ajustar um modelo linear a esses pontos:
\begin{lstlisting}
model2 = LinearRegression(fit_intercept=False)
x = np.vstack([sample1, sample2, sample3])
model2.fit(x, y)
print(model2.coef_, model2.intercept_)
\end{lstlisting}

Vamos imprimir os vários componentes envolvidos.
\begin{lstlisting}
b = model2.coef_[:, np.newaxis]
print("x:\n", x)
print("b:\n", b)
print("y:\n", y[:, np.newaxis])
print("product:\n", np.matmul(x, b))
\end{lstlisting}

\section{Regressão Polinomial}

Pode ser uma surpresa que é possível ajustar uma curva polinomial aos pontos de dados utilizando regressão linear. O truque é adicionar novas variáveis explicativas ao modelo. Abaixo, temos uma única feature $x$ com valores $y$ associados, dados por um polinômio de terceiro grau, com algum ruído (gaussiano) adicionado. Torna-se claro que não podemos explicar os dados bem com uma função linear, então adicionamos duas novas features: $x^2$ e $x^3$. O modelo linear encontrará os coeficientes para essas variáveis.

\begin{lstlisting}
x = np.linspace(-50, 150, 50)
y = 0.15 * x**3 - 20 * x**2 + 5 * x - 4 + 5000 * np.random.randn(50)
plt.scatter(x, y, color="black")

model_linear = LinearRegression(fit_intercept=True)
model_squared = LinearRegression(fit_intercept=True)
model_cubic = LinearRegression(fit_intercept=True)

x2 = x**2
x3 = x**3

model_linear.fit(np.vstack([x]).T, y)
model_squared.fit(np.vstack([x, x2]).T, y)
model_cubic.fit(np.vstack([x, x2, x3]).T, y)

xf = np.linspace(-50, 150, 50)
yf_linear = model_linear.predict(np.vstack([xf]).T)
yf_squared = model_squared.predict(np.vstack([xf, xf**2]).T)
yf_cubic = model_cubic.predict(np.vstack([xf, xf**2, xf**3]).T)

plt.plot(xf, yf_linear, label="linear")
plt.plot(xf, yf_squared, label="squared")
plt.plot(xf, yf_cubic, label="cubic")
plt.legend()

print("Coefficients:", model_cubic.coef_)
print("Intercept:", model_cubic.intercept_)
\end{lstlisting}

Variáveis lineares e ao quadrado não são suficientes para explicar os dados, mas a regressão linear consegue ajustar muito bem uma curva polinomial aos pontos de dados quando variáveis cúbicas são incluídas.

\section{Resumo (Semana 5)}
\begin{itemize}
    \item \texttt{pd.concat} e \texttt{pd.merge} podem combinar dois DataFrames, mas a forma de combinação difere. A função \texttt{pd.concat} concatena com base nos índices, enquanto \texttt{pd.merge} combina com base no conteúdo de variáveis comuns.
    \item A opção \texttt{join="outer"} no \texttt{pd.concat} pode criar valores ausentes, mas \texttt{join="inner"} não. O primeiro retorna a união dos índices, e o segundo, a intersecção.
    \item No \texttt{pd.concat}, índices sobrepostos podem: causar um erro, renumeração de índices, ou criar índices hierárquicos.
    \item O "merging" (mesclagem) pode unir elementos de um-para-um, um-para-muitos, ou muitos-para-muitos.
    \item No agrupamento (grouping), um DataFrame é dividido em DataFrames menores. As principais operações nestes grupos são: aggregate, filter, transform (mantém o formato) e apply.
    \item Series indexadas pelo tempo são chamadas de séries temporais (time series).
    \item A regressão linear pode ser usada para descobrir relacionamentos lineares entre variáveis: pode ter mais de uma feature (variável explicativa) e ajustar polinômios ainda é considerado regressão linear.
\end{itemize}

\vspace{2em}
\noindent\textbf{Créditos:} Este conteúdo foi originalmente criado pela \textit{The University of Helsinki MOOC Center} para o curso \textit{Data Analysis with Python}.
\end{document}
